\chapter{Avaliação}
\label{chap:avaliacao}

\section{Arranjo Experimental}
\label{sec:arranjo}

A avaliação do framework foi conduzida em simulação Gazebo Harmonic com um robô
TurtleBot4 Standard no mundo \textit{warehouse}. O ambiente de simulação garante
condições perfeitamente repetíveis dentro do mesmo laboratório --- sem variação de iluminação, temperatura
ou tráfego de outros agentes --- o que permite isolar a variabilidade introduzida
pelos algoritmos de navegação e SLAM da variabilidade ambiental. Experimentos em
hardware real, com sua camada adicional de incerteza, constituem trabalho futuro.

A configuração da pilha de software é:

\begin{itemize}
    \item \textbf{Sistema operacional}: Ubuntu 24.04 LTS (Noble Numbat)
    \item \textbf{ROS~2}: Jazzy Jalisco (LTS)
    \item \textbf{Simulador}: Gazebo Harmonic 8.x
    \item \textbf{Navegação}: Nav2 (pacotes \texttt{ros-jazzy-nav2-*})
    \item \textbf{SLAM}: SLAM Toolbox (modo síncrono)
    \item \textbf{Plataforma}: TurtleBot4 Standard (via \texttt{nav2\_\allowbreak minimal\_\allowbreak tb4\_\allowbreak sim})
\end{itemize}

\section{Questões de Avaliação}
\label{sec:questoes}

Três questões orientam a avaliação:

\begin{description}
    \item[QA1] O framework consegue executar o protocolo record--replay de forma
    automática e sem intervenção manual entre as fases?

    \item[QA2] Qual é o RMSE inter-replay observado em uma campanha controlada
    com dez execuções independentes iniciadas a partir do mesmo estado? Os
    valores ficam abaixo da tolerância padrão do Nav2 (25\,cm)?

    \item[QA3] A consistência temporal entre execuções equivalentes é suficiente
    para justificar comparações de trajetória baseadas em reamostragem temporal
    normalizada?

    \item[QA4] A repetibilidade observada sob o protocolo controlado (QA2, QA3)
    se mantém quando as execuções ocorrem fora desse regime controlado --- sem
    a mesma ação de navegação, pose inicial fixa ou retorno explícito ao ponto
    de partida entre réplicas? Esta questão foi adicionada durante a execução
    do trabalho, motivada pela disponibilidade de três conjuntos de dados
    exploratórios adicionais no repositório do projeto
    (Seção~\ref{sec:res_exp_complementares}), e por isso é tratada como
    evidência diagnóstica, não como conclusão causal principal.
\end{description}

QA1 corresponde à hipótese H1; QA2 e QA3, em conjunto, correspondem a H2; QA4
corresponde a H3 (Seção~\ref{sec:perguntas_pesquisa}). Não há uma questão de
avaliação dedicada a H4 (custo de engenharia): como discutido na
Seção~\ref{sec:retomada_pp}, essa hipótese é avaliada apenas qualitativamente,
não por um critério mensurável como QA1--QA4.

\section{Procedimento}
\label{sec:procedimento}

Para a campanha final de avaliação, o procedimento é:

\begin{enumerate}
    \item Iniciar uma sessão limpa de simulação Gazebo Harmonic com TurtleBot4
    Standard.
    \item Aguardar ativação completa do Nav2, SLAM Toolbox, sensores e serviços
    do orquestrador.
    \item Gravar uma baseline limpa (\texttt{dissertation\_\allowbreak clean01}) por
    navegação sequencial pelos waypoints definidos para a rota curta de validação.
    \item Armazenar a rota resultante em YAML e o bag MCAP da baseline.
    \item Executar dez replays independentes da rota gravada.
    \item Antes de cada replay, relançar completamente a simulação, a pilha Nav2,
    o orquestrador e o coletor de sensores, garantindo que a odometria inicie no
    mesmo estado em todas as execuções.
    \item Em cada replay, aguardar explicitamente o mapa do SLAM, a estabilização
    do costmap e o estado \texttt{active} dos nós lifecycle do Nav2 antes de enviar
    a rota.
    \item Coletar \texttt{/\allowbreak scan}, \texttt{/\allowbreak odom}, \texttt{/\allowbreak imu} e
    \texttt{/\allowbreak pose} em formato MCAP.
    \item Executar \texttt{analyze\_\allowbreak runs.py} sobre os dez bags, usando
    \texttt{/\allowbreak odom} como fonte da trajetória e reamostragem temporal
    normalizada.
\end{enumerate}

\section{Critérios de Validação}
\label{sec:criterios}

O framework é considerado válido se:

\begin{enumerate}
    \item O protocolo completo (baseline, dez replays e análise) for executado
    com orquestração automática das fases principais (\textbf{QA1}).

    \item Todos os replays válidos começarem do mesmo estado inicial medido em
    \texttt{/\allowbreak odom}.

    \item O RMSE médio e máximo contra a execução de referência forem inferiores
    a 25\,cm (\textbf{QA2}).

    \item O erro final médio e máximo permanecerem inferiores a 25\,cm.

    \item O desvio de duração entre execuções equivalentes for inferior a 10\%
    da duração média (\textbf{QA3}).
\end{enumerate}

O limiar de 10\% acima é usado como critério mínimo de aceitação para QA3. Ele
não significa que durações diferentes sejam desejáveis, apenas que variações
abaixo desse patamar são pequenas o suficiente para sustentar a comparação por
tempo normalizado no cenário controlado. A campanha final ficou abaixo desse
limiar, como apresentado na Seção~\ref{sec:res_qa3}.

\section{Ameaças à Validade}
\label{sec:ameacas}

Seguindo a prática de reportar ameaças à validade em pesquisa empírica de
Engenharia de Software \cite{amigoni2010}, e em resposta ao diagnóstico de
rigor metodológico limitado --- comparação estatística fraca, baixa
disponibilidade de artefatos --- identificado em experimentos de robótica
autônoma publicados \cite{amigoni2014}, esta seção examina os limites da
avaliação nas quatro dimensões clássicas, na linha do que
\citeonline{bonsignorio2015} defendem como pesquisa em robótica replicável e
mensurável.

\textbf{Validade de construto.} O RMSE foi calculado a partir do tópico
\texttt{/\allowbreak odom}, não de um sistema de captura de movimento externo. Tanto a
odometria quanto a pose estimada pelo SLAM são gravadas nos bags, o que permite
comparações cruzadas futuras com uma referência externa. Em simulação, a
odometria é gerada pelo plugin DiffDrive do Gazebo e, portanto, é menos ruidosa
que encoders físicos reais; em robôs físicos, slip de roda e condições do piso
exigiriam uma referência externa ou comparação direta com a pose do SLAM.

Um risco de validade de construto foi identificado e corrigido durante o
desenvolvimento: o modo automático de detecção de tópico do
\texttt{analyze\_\allowbreak runs.py} originalmente só reconhecia um tópico chamado
literalmente \texttt{amcl\_\allowbreak pose}; como este trabalho usa SLAM Toolbox, e não
AMCL, esse tópico nunca existe, e a análise recaía silenciosamente em
\texttt{/\allowbreak odom} mesmo quando a pose corrigida do SLAM
(\texttt{/\allowbreak pose}) também havia sido gravada no mesmo bag. Esse
comportamento poderia inflar artificialmente o RMSE em rotas mais longas ou
demoradas, já que a odometria bruta deriva continuamente enquanto a pose do
SLAM é corrigida pelo casamento de \textit{scans} contra o mapa. O modo
automático foi corrigido para priorizar, nessa ordem, \texttt{amcl\_\allowbreak pose}
$>$ \texttt{pose} (SLAM Toolbox) $>$ \texttt{odom}, escolhendo a primeira opção
com mensagens de fato gravadas; a seleção também pode ser forçada
explicitamente. Análises realizadas antes dessa correção que tenham caído no
\textit{fallback} de \texttt{/\allowbreak odom} deveriam, a rigor, ser
reconsideradas caso a pose do SLAM também estivesse disponível no mesmo bag.

Uma segunda fonte de viés de construto, evidenciada em execuções exploratórias
reais (Seção~\ref{sec:res_exp_complementares}), é que a fase de gravação e a
fase de replay não navegam pelo mesmo mecanismo: \texttt{record} envia uma
sequência de comandos \texttt{go\_\allowbreak to\_\allowbreak point}, um objetivo Nav2 de
cada vez, com o robô assentando em cada waypoint antes de o próximo ser
enviado; \texttt{replay} envia uma única \texttt{play\_\allowbreak route}, percorrendo a
rota inteira como navegação contínua. Em uma campanha exploratória com essa
configuração (rota de quatro waypoints, $N=3$ replays), o RMSE par-a-par entre
replays foi da ordem de 0,02\,m, contra cerca de 0,14\,m entre a baseline e
cada replay, com a baseline durando cerca de 2,5$\times$ menos que os replays
--- diferença explicada quase inteiramente pelo mecanismo de navegação
distinto, não por falta de repetibilidade do robô. Por esse motivo, o RMSE
calculado entre replays entre si (que compartilham o mesmo mecanismo de
navegação) é a métrica mais fiel à repetibilidade que este trabalho pretende
medir --- seguindo a mesma lógica metodológica adotada por
\citeonline{maset2022} para repetibilidade em mapeamento 3D, de comparar
execuções nominalmente idênticas diretamente entre si em vez de cada uma
isoladamente contra um único valor de referência; comparações contra a
gravação baseline, que usa um mecanismo de
navegação distinto, misturam essa diferença de implementação com
repetibilidade real e devem ser interpretadas com essa ressalva.

\textbf{Validade interna.} Campanhas preliminares sem reinicialização do estado
entre replays produziram trajetórias cujos pontos iniciais coincidiam com os
pontos finais da execução anterior, misturando erro de repetição com erro de
inicialização. Por esse motivo, tais campanhas foram tratadas como diagnósticas,
e não como resultado final. A campanha reportada nesta dissertação corrige essa
ameaça ao relançar a simulação e os nós de navegação/coleta antes de cada replay,
fazendo com que os dez bags analisados comecem praticamente no mesmo ponto em
\texttt{/\allowbreak odom}. Ainda assim, fatores como variação interna do controlador
local, estado do SLAM e estabilização do costmap continuam sendo possíveis fontes
de variabilidade residual. Em particular, o controlador local usado
(\texttt{nav2\_\allowbreak mppi\_\allowbreak controller}, algoritmo MPPI --- \textit{Model
Predictive Path Integral} \cite{williams2016}) é estocástico por projeto: a
cada ciclo de controle
(20\,Hz), com \texttt{regenerate\_\allowbreak noises: true}, ruído gaussiano é
reamostrado sobre a trajetória candidata. Ou seja, duas execuções da mesma
rota, no mesmo robô e no mesmo ambiente, divergem parcialmente apenas pela
amostragem aleatória do controlador --- isso não é falha de repetibilidade do
sistema avaliado, é ruído esperado e não-eliminável do método de controle
escolhido (a menos que se troque por um controlador determinístico, como DWB
ou Regulated Pure Pursuit, alternativa não avaliada nesta dissertação).
Adicionalmente, o critério de chegada do Nav2 (\texttt{general\_\allowbreak goal\_\allowbreak
checker}) aceita até 0,25\,m de erro de posição e 0,25\,rad (aproximadamente
14°) de erro angular; desvios de endpoint dentro dessa faixa refletem o
comportamento configurado do Nav2, não necessariamente imprecisão do
framework de medição.

\textbf{Validade externa.} A validação usa simulação com TurtleBot4 em um único
ambiente (\textit{warehouse}). Hardware real introduziria slip de roda, ruído de
sensores, erro de calibração e efeitos de superfície não modelados pelo Gazebo.
Ainda assim, a arquitetura foi projetada para minimizar dependências
específicas de plataforma: o orquestrador depende de actions, tópicos e
transforms padrão do ROS~2/Nav2 --- não de interfaces proprietárias do
TurtleBot4 --- o que favorece, em princípio, a generalização para outras
plataformas com a mesma pilha de navegação. Essa generalização não foi
testada empiricamente nesta dissertação com nenhuma outra plataforma.

\textbf{Validade de conclusão.} A campanha final controlada utiliza $N=10$
replays, o que melhora substancialmente a estimativa de variabilidade em relação
às execuções exploratórias iniciais. Ainda assim, os resultados permanecem
restritos a um único cenário, uma única rota curta, um único robô simulado e uma
única configuração de Nav2. Portanto, embora a campanha seja suficiente para
demonstrar a capacidade do framework de automatizar e medir repetibilidade sob
condições controladas, ela não deve ser interpretada como caracterização
estatística geral do comportamento do TurtleBot4, do Nav2 ou do SLAM Toolbox em
qualquer ambiente. A opção é, portanto, pela profundidade de análise por
execução em detrimento da amplitude de cenários: ao contrário de validações
agregadas de longa distância como o Marathon~2 \cite{macenski2020marathon},
que percorrem dezenas de quilômetros em ambientes variados mas reportam
critérios agregados (distância total, intervenções humanas) sem análise de
consistência entre execuções individuais, esta dissertação restringe o
escopo ambiental para poder medir exatamente essa consistência.

\section{Considerações Finais}
\label{sec:aval_conclusao}

O arranjo de avaliação foi concebido para responder às questões de forma direta e
verificável. A escolha por simulação sacrifica o realismo físico em favor do controle
experimental --- uma troca razoável para uma primeira validação de framework,
onde o objetivo é demonstrar que o sistema funciona como projetado, não que o
hardware físico se comporta de determinada forma. Os resultados são apresentados
no capítulo seguinte.
